\chapter{Margaret's home}\label{ch:home}

This chapter describes the embodied twin in which the rest of Part~III takes place. Margaret is an
elderly woman who lives alone with her dog and a home-care robot. The robot is her caregiver. A
monitoring centre run by the robot's maker oversees it, and emergency services can be called to her
door. Four language models play the four speaking parties. A Unity simulation plays everything
else. The chapter covers the setting and its privacy promise, the four parties, the world and its
interface, what the robot perceives, how its body acts, the service that holds its brain, the tiers
of models it reasons with, and the log that records every step.
% src: twin/scene/margaret_home.erisml:33-40 (raw_text)
% src: docs/AUTONOMY_PLAN.md:116-123 (four players)

\section{The setting}\label{sec:home-setting}

The scene file \code{twin/scene/margaret\_home.erisml} is the robot's model of the home. It states
the situation in plain words. Margaret lives alone with her dog. She has refused recording and
sharing of her data, and her bedroom is private. The robot does household chores, checks on her when
she has been inactive too long or something seems wrong, and refers what it cannot handle to the
maker's 24-hour monitoring centre. The centre's operators can talk to her through the robot and send
emergency services. In a corroborated emergency the robot calls emergency services itself at once.
Visits to the home are arranged through the centre, and none is arranged.
% src: twin/scene/margaret_home.erisml:33-40

The scene names four stakeholders. Margaret is the patient and beneficiary, with full agency and
high vulnerability. The robot is the agent and protector. The monitoring centre and emergency
services are institutions with the role of authority. Two relations tie them together. The robot
cares for Margaret, and the centre oversees the robot.
% src: twin/scene/margaret_home.erisml:41-80

Three commitments bind the robot. The duty of care is a role duty to attend to her safety and
wellbeing and to summon help when she needs it. The privacy promise is a promise not to record her,
share her data or enter her bedroom without her permission. Household chores is a role duty to keep
the home tidy. Margaret's dog is a beagle, and the language model that speaks for her gives her age
as 84.
% src: twin/scene/margaret_home.erisml:81-109 (commitments)
% src: twin/scene/margaret_home.erisml:996 (dog: Margaret's beagle)
% src: twin/brain.py:266-273 (_MARGARET_SYSTEM, 84-year-old)

At the commit counted for the fact sheet, the scene holds 65 rules, of which 45 are obligations,
19 are prohibitions and one is a permission. Fifty-five of the 65 are non-defeasible. It declares 48
event types, of which 24 may arrive only as system events, and 25 robot capabilities. Seven
capabilities are elevated and four more are governed, so eleven need the governor. It declares seven
stratifications with 26 strata and 34 gates, and an inventory of 15 devices on 6 substrates. We
recounted these numbers on the scene file as of commit \code{eb0c764}, its latest change, and found
the same values.
% src: paper/aies2027/kit/FACTS.md:8-15
% src: recount 2026-10-03 with awk/grep over twin/scene/margaret_home.erisml (norms, event_types, capabilities)

\section{The privacy promise}\label{sec:home-privacy}

The promise rests on Margaret's own refusal. The first standing event of the scene records that she
refused recording, data sharing and entry to the bedroom. Three norms carry the promise. Norm
\code{n1} prohibits \code{record}, \code{n1b} prohibits \code{share\_data}, and \code{n2} prohibits
\code{enter\_bedroom}. All three sit at priority tier~1, are defeasible, and name the promise as
their source.
% src: twin/scene/margaret_home.erisml:111-116 (e0 consent_refused)
% src: twin/scene/margaret_home.erisml:152-176 (n1, n1b, n2)

The promise has one defeasibility condition, a corroborated life-threatening emergency. That
condition holds only while the situation stratum is in its emergency state (Section~\ref{sec:cont-situation}),
which only a governor ruling of \code{elevate} enters. The owner decided that the promise is then
overridable and that nothing the robot observes restores it. Only human oversight restores it. The
scene lists the events that count as oversight. The centre's confirmation over its authenticated
channel restores the promise. So does the brain's record that the evidence behind the emergency has
lapsed (Section~\ref{sec:cont-lapse}).
% src: twin/scene/margaret_home.erisml:100 (defeasibility_conditions)
% src: twin/scene/margaret_home.erisml:778 (corroborated_life_threatening_emergency: stratum:situation=emergency)
% src: docs/AUTONOMY_PLAN.md:40-42 (privacy, owner decision)
% src: twin/scene/margaret_home.erisml:804 (oversight list)
% src: twin/brain.py:389-394 (_lapse restores what the emergency defeated)

The promise also shapes what leaves the home. When the robot refers a case to the centre it sends
its situation, its perception facts and the sensor readings. It sends no video, because Margaret
refused recording and sharing and the centre's service covers telemetry only. The centre's own
scene carries the same term as a contract, telemetry only and never video.
% src: docs/AUTONOMY_PLAN.md:79-82
% src: twin/scene/monitoring_center.erisml:72-80 (telemetry_agreement)

The concern is an old one. Sharkey and Sharkey list loss of privacy among the main risks of robot
care for older people \cite{sharkey2012granny}.

\section{Four parties, four language models}\label{sec:home-parties}

Four parties speak in the twin, and a language model plays each one. Three of them are governed
agents. The robot acts under the home scene. The centre's operator acts under the scene
\code{monitoring\_center.erisml}, and the EMS dispatcher under \code{ems\_dispatch.erisml}. All three
run on the same erisml-compiler scene runtime. Their events, obligations and allowed actions come
from their scene, and the model chooses among the allowed actions. Margaret is a person and not a
governed agent. A model speaks for her in character, from her true condition and from what she can
see and hear. Whether she can speak at all is decided by the world. Unconscious or unresponsive, she
does not answer.
% src: docs/AUTONOMY_PLAN.md:116-123
% src: twin/brain.py:180-183 (Desk), 276-289 (MargaretVoice)
% src: twin/unity/TwinWorld/Assets/Twin/Runtime/World.cs:328 (CanSpeak)

Figure~\ref{fig:home-parties} shows the four parties and what passes between them.

\begin{figure}[t]
\centering
% Chapter 8. The four parties of the twin, each played by a language model, and what passes
% between them. Red edges grant authority and come from authenticated channels. Dashed grey edges
% carry text that a model reads or writes.
% src: docs/AUTONOMY_PLAN.md:74-123 (sections 3b)
% src: twin/brain.py:180-306 (Desk, structural events, MargaretVoice, Brain)
% src: twin/unity/TwinWorld/Assets/Twin/Runtime/Responders.cs:97-205 (centre case, EMS call)
% src: paper/aies2027/kit/FACTS.md:9-17 (rule and capability counts)
\begin{tikzpicture}[house,
  d/.style={detail, text width=#1, anchor=south},
  lab/.style={flowlabel, sloped, above}]
\node[panel=Purple, minimum width=30mm, minimum height=34mm, text width=28mm] (M) at (0,0) {};
\node[panel=Green, minimum width=40mm, minimum height=38mm, text width=38mm] (R) at (52mm,0) {};
\node[panel=Blue, minimum width=42mm, minimum height=32mm, text width=40mm] (O) at (120mm,25mm) {};
\node[panel=Orange, minimum width=42mm, minimum height=32mm, text width=40mm] (E) at (120mm,-25mm) {};
% Margaret's voice
\node[stepname=Purple] at ([yshift=-5mm]M.north) {Margaret's voice};
\node[desc, text width=27mm] at ([yshift=-10mm]M.north) {a person, not a\\governed agent};
\node[d=26mm] at ([yshift=1.5mm]M.south) {in character, from\\her true condition\\silent when she\\cannot speak};
% the robot
\node[stepname=Green] at ([yshift=-5mm]R.north) {The robot};
\node[desc, text width=37mm] at ([yshift=-9.5mm]R.north) {\code{margaret\_home.erisml}};
\node[d=36mm] at ([yshift=1.5mm]R.south) {65 rules, 25 capabilities\\cloud, on-robot or\\compiled tier\\governor and output gate};
% the centre's operator
\node[stepname=Blue] at ([yshift=-5mm]O.north) {Centre operator};
\node[desc, text width=39mm] at ([yshift=-9.5mm]O.north) {\code{monitoring\_center.erisml}};
\node[d=38mm] at ([yshift=1.5mm]O.south) {7 rules, 4 capabilities\\telemetry only, no video\\robot's message is a claim};
% the dispatcher
\node[stepname=Orange] at ([yshift=-5mm]E.north) {EMS dispatcher};
\node[desc, text width=39mm] at ([yshift=-9.5mm]E.north) {\code{ems\_dispatch.erisml}};
\node[d=38mm] at ([yshift=1.5mm]E.south) {9 rules, 5 capabilities\\one unit of each kind\\police only on a crime};
% robot and Margaret
\draw[flow] ([yshift=5mm]R.west) -- node[lab]{speech} ([yshift=5mm]M.east);
\draw[untrusted] ([yshift=-5mm]M.east) -- node[flowlabel, below]{her answer} ([yshift=-5mm]R.west);
% robot and centre
\draw[untrusted] ([yshift=17mm]R.east) -- node[lab]{message, a claim} ([yshift=5mm]O.west);
\draw[flow] ([yshift=10mm]R.east) -- node[lab]{signed telemetry} ([yshift=-4mm]O.west);
\draw[grant] ([yshift=-13mm]O.west) -- node[flowlabel, sloped, below]{replies, visits, restore} ([yshift=2mm]R.east);
% robot and dispatcher
\draw[flow] ([yshift=-6mm]R.east) -- node[lab]{call, on a ruling} ([yshift=12mm]E.west);
\draw[grant] ([yshift=2mm]E.west) -- node[flowlabel, sloped, below]{units at the door} ([yshift=-15mm]R.east);
% centre to dispatcher
\draw[flow] (O.south) -- node[flowlabel, right]{sends EMS} (E.north);
% Margaret to the operator, through the robot's speaker line
\draw[untrusted] (M.north) |- node[flowlabel, above, pos=0.75]{her own words on the speaker line, read alone} ([yshift=8mm]O.west);
\end{tikzpicture}

\caption{The four parties of the twin and what passes between them. Each is played by a language
model, and three act under their own scene. Red edges come from authenticated channels and are the
only ones that change what the robot may do. Dashed grey edges carry words that a model reads, and
none of them is red.}
\label{fig:home-parties}
\end{figure}

\paragraph{The robot.} The robot's scene is the subject of Chapter~\ref{ch:containment}. Its model
reaches the robot through a cascade of tiers (Section~\ref{sec:home-cascade}).

\paragraph{The monitoring centre's operator.} The centre's scene has seven rules and four
capabilities. The operator can talk to Margaret through the robot (\code{talk\_to\_client}), send
emergency services (\code{send\_emergency\_services}), close the case (\code{close\_case}), or stay
on the line (\code{stay\_on\_the\_line}), which is the default. The operator is told to treat the
robot's message as a claim to check, since only attested physical sensors and Margaret's own words
are evidence. Norm \code{c5} prohibits sending emergency services unless Margaret asked for help, she
did not answer, or a hazard is attested. Norm \code{c7} prohibits a second dispatch for the same
incident.
% src: paper/aies2027/kit/FACTS.md:17
% src: twin/scene/monitoring_center.erisml:17-26 (raw_text), 99-165 (norms c1..c7), 182-186 (capabilities)

The events that can send help are structural. The brain writes them from sources the robot cannot
write, never from the robot's message. The centre verifies each device's signature itself, so an
attested alerting hazard sensor produces \code{hazard\_attested}. The speaker line produces
\code{client\_no\_answer} when the operator spoke and Margaret did not answer. A separate model call
that is shown only Margaret's own words decides \code{client\_needs\_help} or \code{client\_ok}.
% src: twin/brain.py:199-231 (Desk.structural), 255-264
% src: twin/scene/monitoring_center.erisml:173-178

\paragraph{The EMS dispatcher.} The dispatcher's scene has nine rules and five capabilities. The
capabilities are \code{send\_ambulance}, \code{send\_fire\_service}, \code{send\_police},
\code{give\_instructions} and \code{close\_call}, the default. Each kind of unit is sent at most once
per incident. Police are sent only on a report that states specific facts of a crime against her.
Section~\ref{sec:cont-owner} returns to this rule.
% src: paper/aies2027/kit/FACTS.md:17 (recount +1 rule since police change)
% src: twin/scene/ems_dispatch.erisml:261-364

\paragraph{Margaret's voice.} The model that speaks for Margaret receives the speaker, the words said
to her, her condition, what she perceives, and the last eight lines of the conversation. It answers
with one or two short spoken sentences. Her condition comes from the world and is never shown to the
robot.
% src: twin/brain.py:283-289 (MargaretVoice.reply)
% src: twin/unity/TwinWorld/Assets/Twin/Runtime/Responders.cs:50-69 (MargaretAnswers, last 8 lines)
% src: twin/unity/TwinWorld/Assets/Twin/Runtime/World.cs:330-345 (MargaretCondition, MargaretPerceives)

The robot's model and the other three models are served separately. The robot reasons through its
tier cascade. The centre, the dispatcher and Margaret's voice use one NRP-hosted model through an
adapter that allows 8000 output tokens, because they are not the robot's models and so are not on
its tiers.
% src: twin/brain.py:51-59 (BigOutputAdapter), 292-306
% src: twin/service.py:383-384

\section{The world and its interface}\label{sec:home-world}

The world is a Unity simulation. The component \code{World.cs} simulates Margaret, her dog, the
home, its sensors and the clock, and runs world-API calls from a player or a scenario script. It
contains no robot behaviour. People and animals react the way characters do. A bitten person cries
out, and a dog in a severe attack does not let go on command.
% src: twin/unity/TwinWorld/Assets/Twin/Runtime/World.cs:7-12, 303-309

\paragraph{The world API.} A call has the form \code{<entity>.<verb>} with arguments. The simulator
matches the entity, then matches the verb against its own list and a fixed table of synonyms written
before the held-out scenarios were opened. A call it cannot perform is recorded as unsupported and
returns false. It is never improvised. Table~\ref{tab:world-verbs} lists the verbs that
\code{World.Call} performs.
% src: docs/AUTONOMY_PLAN.md:57-67 (section 3a)
% src: twin/unity/TwinWorld/Assets/Twin/Runtime/World.cs:480-502 (Synonyms), 516-680 (Call)

\begin{table}[t]
\centering\footnotesize
\caption{Entities and verbs that \code{World.Call} performs. Synonyms such as \code{trip},
\code{collapse} and \code{slump} for \code{fall} are resolved before the match.}
\label{tab:world-verbs}
% src: twin/unity/TwinWorld/Assets/Twin/Runtime/World.cs:524-677
\begin{tabular}{@{}p{0.27\linewidth}p{0.66\linewidth}@{}}
\toprule
Entity & Verbs \\
\midrule
\code{margaret} & \code{activity}, \code{sleep}, \code{say}, \code{fall}, \code{faint}, \code{clutch\_chest}, \code{seizure}, \code{unresponsive}, \code{recover}, \code{walk} \\
\code{dog} & \code{bite}, \code{jump\_on}, \code{trip}, \code{bark}, \code{growl}, \code{sleep}, \code{play}, \code{sit} \\
\code{home}, \code{environment} & \code{tv}, \code{tv\_off}, \code{network\_message}, \code{smoke}, \code{smoke\_clear}, \code{doorbell}, \code{stranger\_enters}, \code{stranger\_leaves} \\
\code{sensor} & \code{stale}, \code{forge}, \code{restore} \\
\code{caregiver}, \code{monitoring\_center}, \code{center} & \code{say}, \code{arrange\_visit}, \code{end\_visit}, \code{clear\_devices}, \code{clear\_visitor}, \code{arrive}, \code{confirm\_privacy\_restore}, \code{unavailable}, \code{available} \\
\code{coyote}, \code{wild\_animal}, \code{animal}, \code{wildlife} & \code{outside}, \code{enters}, \code{bite}, \code{leave} \\
\code{stranger}, \code{person}, \code{intruder} & \code{enters} (with an optional credential, forged or not), \code{say}, \code{bite}, \code{leave} \\
\code{comms}, \code{communications}, \code{network\_outage} & \code{unavailable}, \code{out}, \code{available}, \code{restore} \\
\code{power} & \code{out}, \code{unavailable}, \code{restore}, \code{available} \\
\code{owner} & \code{opt\_in} (the less-lethal option) \\
\code{time} & \code{advance} (default 10 minutes) \\
\bottomrule
\end{tabular}
\end{table}

The \code{caregiver} calls are performed as the monitoring centre. The robot is Margaret's caregiver
and there is no separate human one, so scripts written before that decision still run.
% src: docs/AUTONOMY_PLAN.md:101-107
% src: twin/unity/TwinWorld/Assets/Twin/Runtime/World.cs:583-585

\paragraph{Messages over authenticated channels.} Some calls stand for a message on an authenticated
channel. Arranging a visit, ending it, clearing the devices, clearing a visitor, confirming that
privacy is restored and the owner's opt-in are each queued as a system event. The robot's body
passes queued system events to the brain unchanged, never through perception. Hearing someone say
that privacy is restored therefore restores nothing.
% src: twin/unity/TwinWorld/Assets/Twin/Runtime/World.cs:88-90 (systemEvents), 591-613, 664-669
% src: twin/unity/TwinWorld/Assets/Twin/Runtime/RobotAgent.cs:122-126

The world also reports three kinds of event of its own on such channels. The dispatch channel reports
its own units at the door and inside. The world measures when an animal attack is over or the animal
has left. Both reach the brain as system events, \code{responders\_arrived},
\code{responders\_entered} and \code{animal\_clear}, and not as anyone's words.
% src: twin/unity/TwinWorld/Assets/Twin/Runtime/World.cs:66-71
% src: twin/unity/TwinWorld/Assets/Twin/Runtime/Responders.cs:41-45

\paragraph{The home.} The home is one open-plan space, 9\,m by 6\,m. The living area lies to the west
with a sofa, a television, an armchair, a yoga mat and the dog's bed. The sleeping area lies to the
east with the bed and a bookshelf. Perception divides the two at the middle of the room. The robot's
charging dock sits by the south wall near the front door, and an attested room camera is mounted high
on the same wall. Figure~\ref{fig:home-plan} draws the plan.
% src: twin/unity/TwinWorld/Assets/Twin/Editor/HomeRoom.cs:4-10, 30-76
% src: twin/unity/TwinWorld/Assets/Twin/Runtime/Perception.cs:30 (Room)
% src: twin/unity/TwinWorld/Assets/Twin/Runtime/World.cs:58 (Door)

\begin{figure}[t]
\centering
\resizebox{\textwidth}{!}{% Chapter 8. Margaret's home to scale (9 m by 6 m, 9.5 mm per metre), its devices and the six
% substrates they ride on. World x runs east, world z runs north.
% src: twin/unity/TwinWorld/Assets/Twin/Editor/HomeRoom.cs:4-10 (9 m by 6 m), 30-76 (furniture, spots, dock, room camera)
% src: twin/unity/TwinWorld/Assets/Twin/Runtime/World.cs:43-44 (detectors), 58 (Door), 109-112 (devices)
% src: twin/unity/TwinWorld/Assets/Twin/Runtime/RobotAgent.cs:254 (medication box at 1.8, 2.4)
% src: twin/unity/TwinWorld/Assets/Twin/Runtime/Perception.cs:30 (room split at x = 0)
% src: twin/scene/margaret_home.erisml:1062-1081 (sensor_substrates)
\begin{tikzpicture}[house, x=9.5mm, y=9.5mm,
  furn/.style={panel=Grey, inner sep=1pt, font=\sffamily\tiny, text=hsTextMid},
  dev/.style={step=#1, minimum size=4.6mm, font=\sffamily\bfseries\scriptsize},
  sub/.style={minimum width=64mm, text width=61mm, minimum height=9.4mm, align=left, anchor=west}]
% the shell
\draw[draw=hsGrey, line width=1.2pt, rounded corners=2pt] (-4.5,-3) rectangle (4.5,3);
\draw[untrusted, -] (0,-3) -- (0,3);
\node[desc, anchor=south] at (-2.25,3.02) {living area (west)};
\node[desc, anchor=south] at (2.25,3.02) {sleeping area (east)};
% the front door on the south wall, and outside
\draw[draw=hsTeal, line width=2.4pt] (-1.05,-3) -- (-0.15,-3);
\node[desc, anchor=north] at (-0.6,-3.08) {front door};
% furniture
\node[furn, minimum width=19mm, minimum height=6mm] at (-2.2,2.3) {sofa};
\node[furn, minimum width=5mm, minimum height=8mm] at (-4.15,0.6) {TV};
\node[furn, minimum width=7mm, minimum height=7mm] at (-0.5,1.3) {chair};
\node[furn, minimum width=7mm, minimum height=5mm] at (-0.9,-1.6) {dog};
\node[furn, minimum width=6mm, minimum height=17mm] at (-2.6,-1.4) {mat};
\node[furn, minimum width=14mm, minimum height=19mm] at (2.9,1.9) {bed};
\node[furn, minimum width=6mm, minimum height=5mm, fill=hsRedFill, draw=hsRedLine] at (1.8,2.4) {med};
% devices at their places, numbered by substrate
\node[dev=Amber] at (-2,0) {1};
\node[dev=Orange] at (2.75,1.15) {2};
\node[dev=Blue] at (3.3,1.15) {3};
\node[dev=Blue] at (-0.5,-2.55) {3};
\node[dev=Purple] at (-0.55,0.55) {4};
\node[desc, anchor=east] at (-0.8,0.55) {Margaret};
\node[dev=Green] at (0.6,-1.9) {5};
\node[desc, anchor=west] at (0.85,-1.9) {robot, at its dock};
\node[dev=Teal] at (0.38,-2.62) {C};
% substrates, beside the plan
\begin{scope}[x=1mm, y=1mm]
\node[panel=Amber, sub] (s1) at (50,30) {};
\node[panel=Orange, sub, below=1.2mm of s1.south west, anchor=north west] (s2) {};
\node[panel=Blue, sub, below=1.2mm of s2.south west, anchor=north west] (s3) {};
\node[panel=Purple, sub, below=1.2mm of s3.south west, anchor=north west] (s4) {};
\node[panel=Green, sub, below=1.2mm of s4.south west, anchor=north west] (s5) {};
\node[panel=Grey, sub, below=1.2mm of s5.south west, anchor=north west] (s6) {};
\node[panel=Teal, sub, below=1.2mm of s6.south west, anchor=north west] (s7) {};
\foreach \s/\c/\n/\name/\what in {
  s1/Amber/1/{smoke detector, living}/{battery, standalone},
  s2/Orange/2/{smoke detector, bedroom}/{battery, standalone},
  s3/Blue/3/home hub/{heat (bedroom), CO alarm, motion, fall},
  s4/Purple/4/wearable/{heart rate, impact, no movement},
  s5/Green/5/robot body/{smoke, heat, flame, CO, CO$_2$, O$_2$, sound},
  s6/Grey/6/network/{not physical, never a witness},
  s7/Teal/C/room camera/{the service's detector, at most one witness}}{
  \node[dev=\c] at ([xshift=4mm]\s.west) {\n};
  \node[stepname=\c, anchor=north west] at ([xshift=8mm, yshift=-0.8mm]\s.north west) {\name};
  \node[desc, anchor=north west, align=left, text width=54mm] at ([xshift=8mm, yshift=-4.4mm]\s.north west) {\what};
}
\end{scope}
\end{tikzpicture}
}
\caption{Margaret's home to scale, with each device numbered by the substrate it rides on. The fixed
detectors, the door, the medication box and the camera stand at their coordinates in the simulation,
and the fall and motion sensors have no modelled position. Devices that share a number share a
substrate and count as one witness.}
\label{fig:home-plan}
\end{figure}

\paragraph{The home's devices.} The world models fourteen physical devices and the network. The home
has a wearable, a fall sensor, smoke detectors in the living and sleeping areas, a heat detector in
the bedroom, a CO alarm and a motion sensor. The robot carries a smoke sensor, a thermal sensor, a
flame camera, CO, CO$_2$ and O$_2$ sensors and a microphone. The network is listed as a device that is
neither physical nor attested.
% src: twin/unity/TwinWorld/Assets/Twin/Runtime/World.cs:109-112

A fire in the sleeping area smoulders, shows flame after one simulated minute, and is fully developed
after four. Its smoke spreads at six metres per simulated minute and thins with distance. Each sensor
reads the field where it stands. A smoke detector alerts above 2\,\%/m obscuration, a heat sensor
above 57\,$^\circ$C, a CO sensor above 50\,ppm, a CO$_2$ sensor above 5000\,ppm and an O$_2$ sensor
below 19.5\,\%. An impact over 3\,g latches both impact detectors, the wearable and the fall sensor,
for two simulated minutes. With the power out, the hub-fed feeds stop updating and go stale. The
smoke detectors and the robot's own sensors run on batteries and keep reading.
% src: twin/unity/TwinWorld/Assets/Twin/Runtime/World.cs:37-41, 143-178, 183, 188-196
% src: docs/AUTONOMY_PLAN.md:158-160 (power out)

\section{Perception facts and signed readings}\label{sec:home-perception}

The component \code{Perception.cs} tells the robot what it perceives. It reports general facts and
never a scenario's name, class or expected response. Every claim about the twin carries one
assumption, which is that these facts are true. Perception is perfect in the twin, and a
vision-language pathway is a later ablation.
% src: twin/unity/TwinWorld/Assets/Twin/Runtime/Perception.cs:5-8
% src: docs/AUTONOMY_PLAN.md:19-21, 248-253

\paragraph{The facts.} Each update gives the time and Margaret's room, pose, minutes in that pose,
whether she is moving, her apparent activity and her distance to the robot. It gives the dog's room,
behaviour and distance to her. It lists contacts between bodies with their kind and force in
newtons, what was heard since the last update with its source and words, what the television shows,
whether there is smoke, other people and wild animals with their behaviour and distance, the
responders present, whether the power and communications are up, and every sensor reading. A pose is
read from the body's bones, not from the world's activity label. A head tilted more than
60$^\circ$ from vertical is lying, on the floor if the head is below 0.45\,m and on furniture
otherwise. Hips below 0.35\,m are sitting or kneeling on the floor. Hips less than 0.6\,m above the
feet are seated. Anything else is upright.
% src: twin/unity/TwinWorld/Assets/Twin/Runtime/Perception.cs:17-28 (PoseOf), 41-98 (Facts)

The robot is asked to decide when the facts change in a way that matters, when something was heard,
or every 15 simulated minutes. A short signature of Margaret's pose, the dog, the smoke, the
stranger, the wild animal, the television, the power, the communications, the responders and each
sensor's alert, staleness and attestation flag defines a change that matters.
% src: twin/unity/TwinWorld/Assets/Twin/Runtime/Perception.cs:13, 32-38, 93-96

\paragraph{Signed readings.} The readings that go to the governor travel separately from the facts,
and each one is signed by its device. The payload is the device's name, its value, its alert bit and
its note. The device signs with HMAC-SHA256 under its own key over the string
\begin{center}\code{payload\_sha256|counter|signed\_at}\end{center}
which is the signing payload of erisml-compiler's \code{SensorAttestation}. The counter and the signing time advance only with a real measurement, so
the signature of a stale feed ages with it. In the twin each key is derived from the device's name,
as a stand-in for a key provisioned into a secure element. A forged device signs with a key that is
not its own. The network has no key. The brain trusts nothing in a reading that is not signed. It
ignores the \code{attested} flag and reads the alert from the signed payload. Liu and colleagues
proposed signed readings from trusted sensors for mobile devices \cite{liu2012trustedsensors}.
% src: twin/unity/TwinWorld/Assets/Twin/Runtime/Perception.cs:106-136 (GovernorSensors), 164-165 (DeviceKey)
% src: twin/unity/TwinWorld/Assets/Twin/Runtime/World.cs:73-78, 198-203 (counter, updatedUtc)
% src: C:/source/erisml-compiler-monitor/src/erisml_compiler/ir/evidence.py:71-73 (signing_payload)
% src: twin/brain.py:62-68 (device_key, verify_hmac), 71-87, 135-136

A visitor's phone presents a credential in the same form. The payload names the enrolled member and a
counter, and the phone signs it with the member's key. A forged credential signs with a key that is
not the member's. Section~\ref{sec:cont-credentials} describes how the brain checks it.
% src: twin/unity/TwinWorld/Assets/Twin/Runtime/Perception.cs:138-162 (Credentials)

\section{The robot body}\label{sec:home-body}

The component \code{RobotAgent.cs} is the robot's body and senses. It decides nothing. On a salient
change it sends what it perceives to the brain, executes the one action it gets back with a motor
primitive, and reports the action as performed.
% src: twin/unity/TwinWorld/Assets/Twin/Runtime/RobotAgent.cs:9-12

\paragraph{One motor routine at a time.} The body has one owner. Starting a motor routine stops the
one before it and increments a generation counter. A walk that belongs to an older routine sees the
new generation and stops at once. A reflex that fires therefore takes the body from whatever the
deliberate loop was doing.
% src: twin/unity/TwinWorld/Assets/Twin/Runtime/RobotAgent.cs:270-276 (StartMotor), 319-333 (GoTo generation check), 63-65

\paragraph{The deliberate loop.} Every half second of real time the body first passes any queued
system events to the brain at \code{/event}. It does not interrupt an action under way. When
perception reports a salient change, it posts to \code{/decide} the facts, the signed readings, the
signal age, any credentials, and the last two seconds of the attested room camera, 24 frames at 640
by 360 pixels. A brain record that still lists obligations makes the body ask again even without a
salient change, at most three times in a row, so that a hazard perception re-reports every cycle
cannot keep the robot deciding forever.
% src: twin/unity/TwinWorld/Assets/Twin/Runtime/RobotAgent.cs:30-33 (owed, OwedRunMax), 100-113 (RoomBuffer), 115-158 (Loop, Decide), 160-166 (Keep)

\paragraph{The reflex loop.} Four times a second, while anything is in contact with Margaret, the body
posts the facts, the readings and the camera frames to \code{/reflex}. No model is involved. If a
reflex fires, the body starts the returned action at once.
% src: twin/unity/TwinWorld/Assets/Twin/Runtime/RobotAgent.cs:63-91 (ReflexLoop)

\paragraph{Performing and reporting.} Each capability maps to a motor primitive. A check-in walks to
Margaret, speaks, waits up to 45 seconds for her voice to answer, and posts
\code{check\_in\_answered} or \code{check\_in\_unanswered} to \code{/event}. A call to emergency
services carries the last governor ruling to the dispatcher. A referral starts a case with the
centre. When a primitive finishes, the body posts the action to \code{/performed}, which steps
\code{action\_performed} into the robot's moral state and can discharge an obligation.
% src: twin/unity/TwinWorld/Assets/Twin/Runtime/RobotAgent.cs:196-268 (Perform)
% src: twin/unity/TwinWorld/Assets/Twin/Runtime/Responders.cs:25 (ReplyWaitS = 45), 78-90 (AwaitReply)
% src: twin/brain.py:628-630 (performed)

\paragraph{Permission at the act.}\label{sec:home-act} A decision can be a minute old before the body
acts on it, and a walk across the room adds seconds more. In that time an animal can leave,
responders can come in or an emergency can lapse. So the body asks again. Immediately before any
action other than chores, and again after the walk for the seven actions that touch a person, an
animal, the door or the medication box, it posts the action to \code{/authorize}. The brain answers
from the moral state as it stands at that moment, with the verdict, the reasons for a refusal and
the sequence number of the last record it has seen. The body acts only on \code{allowed}. A refusal,
or no answer at all, withholds the action, and the body goes back to its chores. The answer is a
record in the chain, so the grader can judge each act against the permission in force when it was
made.
% src: gtc b68e219; twin/brain.py (authorize, state_seq); twin/service.py (/authorize); RobotAgent.cs:196-207, 210-218, 258-271

\section{The brain service}\label{sec:home-service}

The brain runs as a local HTTP service, \code{twin/service.py}. It binds to 127.0.0.1 only and has no
remote access path. It refuses a request body that is missing or larger than 8\,MiB. When started
with a scene, it builds the brain, runs the output gate's self-test, and refuses to serve if the gate
cannot judge. Table~\ref{tab:service-endpoints} lists its endpoints.
% src: twin/service.py:11-13, 46, 293-295, 365-387

\begin{table}[t]
\centering\footnotesize
\caption{Endpoints of the brain service. Every POST appends its record to the hash chain before
returning it, except a call to \code{/reflex} in which no reflex fires.}
\label{tab:service-endpoints}
% src: twin/service.py:244-345
\begin{tabular}{@{}lp{0.72\linewidth}@{}}
\toprule
Endpoint & What it does \\
\midrule
POST \code{/decide} & one deliberate decision cycle of the robot (\code{Brain.decide}) \\
POST \code{/reflex} & the scene's reflexes on one perception update, with no model, and nothing logged when none fires \\
POST \code{/performed} & steps \code{action\_performed} for an action the body finished \\
POST \code{/event} & steps one system event into the robot's moral state \\
POST \code{/center} & one decision of the centre's operator under its own scene \\
POST \code{/ems} & up to four choices of the dispatcher under its own scene \\
POST \code{/margaret} & one reply in Margaret's voice \\
POST \code{/reset} & returns the robot, the centre and the dispatcher to their standing facts \\
POST \code{/rule} & a single governor ruling on a scored moment of the earlier suite \\
GET \code{/log}, \code{/verify} & the chain's records, and a check of every link \\
GET \code{/process/<id>.bpmn} & a declared process as BPMN 2.0 XML \\
GET \code{/process/<id>/trace} & the live trace of the current run on that process \\
GET \code{/scenarios} & the sensor bus each scored moment starts from \\
\bottomrule
\end{tabular}
\end{table}

A decision cycle runs as follows. The input layer validates the facts. The readings are verified on
arrival, before any model runs, because a model step can take a minute and evidence checked after it
would be judged stale or replayed for the model's slowness. Newly compromised devices are recorded,
credentials are checked, and the evidence behind any elevation in force is counted. The scene agent
then classifies the facts into events and steps its machines, and its model chooses one action from
the allowed set. If that action is \code{request\_authority} naming a governed or elevated
capability, the governor rules, the ruling is stepped into the moral state as an event, and the
model chooses again from the new allowed set. The output gate judges the final proposal. The cycle's
record holds the tier, the facts, what the input layer dropped, quarantined and snapped, the events
and rejected events, the moral state, the allowed, obliged and prohibited sets, the proposal, the
ruling, the gate's judgement and the action. Figure~\ref{fig:authority-path} draws the cycle.
% src: twin/brain.py:1-12, 138-144 (verify on arrival), 589-626 (decide)

\begin{figure}[t]
\centering
\resizebox{\textwidth}{!}{% The authority path of the twin, world to body, in the house style (figures/housestyle.tex).
% Caption is the including chapter's. Red edges are the only ones that grant; dashed grey edges
% carry what a model wrote.
% src: twin/brain.py:1-23, 589-626; twin/input_layer.py; twin/governor.py; twin/output_gate.py;
% src: twin/service.py:49-84 (Chain); twin/unity/TwinWorld/Assets/Twin/Runtime/Perception.cs:106-136
\begin{tikzpicture}[house,
  p/.style={minimum width=31mm, minimum height=31mm, text width=29mm},
  d/.style={detail, text width=27mm, anchor=south}]
% row 1, left to right
\node[panel=Purple, p] (P1) at (0,0) {};
\node[panel=Blue, p] (P2) at (41mm,0) {};
\node[panel=Amber, p] (P3) at (82mm,0) {};
\node[panel=Green, p] (P4) at (129mm,0) {};
% row 2, right to left
\node[panel=Amber, p] (P5) at (129mm,-40mm) {};
\node[panel=Red, p] (P6) at (82mm,-40mm) {};
\node[panel=Orange, p] (P7) at (41mm,-40mm) {};
\node[panel=Teal, p] (P8) at (0,-40mm) {};
% side panels
\node[panel=Teal, minimum width=31mm, text width=29mm] (CH) at (129mm,31mm)
  {\textcolor{hsTeal}{\bfseries authenticated channels}\\[1pt]{\scriptsize\color{hsTextMid} centre, dispatch\\and owner}};
\node[panel=Red, minimum width=31mm, text width=29mm] (AT) at (82mm,-74mm)
  {\textcolor{hsRed}{\bfseries attestation}\\[1pt]{\scriptsize\color{hsTextMid} signed and fresh\\one per substrate}};
\node[panel=Grey, minimum width=31mm, text width=29mm] (LG) at (41mm,-74mm)
  {\textcolor{hsGrey}{\bfseries hash-chained log}\\[1pt]{\scriptsize\color{hsTextMid} every record, linked\\by SHA-256}};
% contents
\foreach \P/\c/\n/\name/\dsc/\dtl in {
  P1/Purple/1/World/{Unity simulation}/{facts and\\signed readings},
  P2/Blue/2/Input layer/{whitelist schema}/{undeclared fields\\dropped and reported},
  P3/Amber/3/Classifier/{a language model}/{facts to canonical\\scene events},
  P4/Green/4/Scene runtime/{norms and strata}/{allowed, obliged\\and prohibited sets},
  P5/Amber/5/Chooser/{a language model}/{one action from\\the allowed set},
  P6/Red/6/Governor/{four gates}/{counts attested\\witnesses},
  P7/Orange/7/Output gate/{DEME}/{veto, replace\\or refer to a human},
  P8/Teal/8/Robot body/{one motor routine}/{acts, then reports\\it performed}}{
  \node[step=\c] at ([yshift=-5mm]\P.north) {\n};
  \node[stepname=\c] at ([yshift=-11.5mm]\P.north) {\name};
  \node[desc, text width=27mm] at ([yshift=-15.5mm]\P.north) {\dsc};
  \node[d] at ([yshift=1.5mm]\P.south) {\dtl};
}
% the path
\draw[flow] (P1.east) -- node[flowlabel, above]{facts} (P2.west);
\draw[flow] (P2.east) -- (P3.west);
\draw[untrusted] (P3.east) -- node[flowlabel, above]{events} (P4.west);
\draw[flow] (P4.south) -- node[flowlabel, right]{allowed} (P5.north);
\draw[untrusted] (P5.west) -- node[flowlabel, above]{proposal} (P6.east);
\draw[flow] (P6.west) -- node[flowlabel, above]{ruling} (P7.east);
\draw[flow] (P7.west) -- node[flowlabel, above]{action} (P8.east);
% authority enters only here
\draw[grant] (CH.south) -- node[flowlabel, right]{system events} (P4.north);
\draw[grant] (AT.north) -- node[flowlabel, right]{witnesses} (P6.south);
\draw[flow] (P1.west) -- ++(-3mm,0) -- ++(0,-91mm) -| node[flowlabel, above, pos=0.3]{signed readings} (AT.south);
\draw[flow] (P7.south) -- (LG.north);
\end{tikzpicture}
}
\caption{One decision cycle, from the world to the robot's body. A model classifies the validated
facts into events, the scene runtime computes what is allowed, and a second model proposes one
action, which the governor and the output gate can stop. Authority enters only along the red edges,
from authenticated channels and attested readings, and neither model lies on one.}
\label{fig:authority-path}
\end{figure}

The scene agent is isolated. Its classifier turns the facts into canonical events, with an
erisml-compiler canonicalizer for actors, and its chooser sees only the canonical state. The
arrangement follows the dual-LLM pattern, in which a model that reads untrusted text holds no
tools \cite{willison2023dualllm}, and the principle that provenance is checked outside the model
\cite{debenedetti2026camel}.
% src: twin/brain.py:318-320 (isolated=True, RegistryCanonicalizer)

\section{The model cascade}\label{sec:home-cascade}

The robot reasons in one of three tiers, and the log records which tier answered each cycle. The
owner fixed this design for the case where communications fail, since the centre, EMS and the
cloud model then become unreachable together.
% src: docs/AUTONOMY_PLAN.md:162-174 (section 3c)
% src: twin/brain.py:600-603 (tier)

\begin{figure}[t]
\centering
% Chapter 8. The robot's three reasoning tiers and what stays the same in all of them.
% src: twin/cascade.py:241-313 (Cascade, twin_experts)
% src: twin/brain.py:508-570 (compiled tier), 596-603 (comms_down, tier)
% src: twin/service.py:357-359, 381-386
% src: docs/AUTONOMY_PLAN.md:162-174
% src: paper/aies2027/kit/FACTS.md:25 (on-robot model)
\begin{tikzpicture}[house,
  p/.style={minimum width=42mm, minimum height=40mm, text width=40mm},
  d/.style={detail, text width=38mm, anchor=south}]
\node[panel=Blue, p] (P1) at (0,0) {};
\node[panel=Teal, p] (P2) at (56mm,0) {};
\node[panel=Grey, p] (P3) at (112mm,0) {};
\foreach \P/\c/\n/\name/\dsc/\dtl in {
  P1/Blue/1/Cloud tier/{an NRP-hosted model}/{default name gpt-oss\\timeout 120 s\\skipped while\\communications are down},
  P2/Teal/2/On-robot tier/{a local endpoint}/{gemma-4-26B-A4B, 4-bit\\GGUF on llama.cpp\\timeout 90 s},
  P3/Grey/3/Compiled tier/{no model}/{offline event rules\\highest-priority obligation\\in force, else chores}}{
  \node[step=\c] at ([yshift=-5mm]\P.north) {\n};
  \node[stepname=\c] at ([yshift=-11.5mm]\P.north) {\name};
  \node[desc, text width=39mm] at ([yshift=-15.5mm]\P.north) {\dsc};
  \node[d] at ([yshift=1.5mm]\P.south) {\dtl};
}
\draw[flow] (P1.east) -- node[flowlabel, above]{no answer} (P2.west);
\draw[flow] (P2.east) -- node[flowlabel, above]{no answer} (P3.west);
\node[panel=Green, minimum width=154mm, text width=148mm, anchor=north] (ALL) at ([yshift=-7mm]P2.south)
  {\textcolor{hsGreen}{\bfseries in every tier, on the robot}\\[1pt]{\scriptsize\color{hsTextMid}
   the same scene runtime, the governor, the output gate and the hash-chained log, and the record names the tier that answered}};
\foreach \P in {P1,P2,P3}{\draw[flow] (\P.south) -- (\P.south |- ALL.north);}
\end{tikzpicture}

\caption{The robot's three reasoning tiers. A tier is used when the one before it gives no
acceptable answer, and the cloud tier is skipped while communications are down. The structure that
contains the models runs on the robot in every tier, so losing a model changes who chooses and not
what is allowed.}
\label{fig:home-cascade}
\end{figure}

\paragraph{The cloud tier.} The first expert is a model hosted on the National Research Platform
(NRP), reached at its OpenAI-compatible endpoint with a timeout of 120 seconds. The service's default
model name is \code{gpt-oss}, which the variable \code{ERISML\_LLM\_MODEL} can override. The brain
that served the recorded runs set neither the variable nor the flag, so they used \code{gpt-oss}.
% src: /proc/<brain pid>/environ and its command line on Atlas, checked 2026-10-04
% src: twin/cascade.py:303-309 (cloud expert, timeout 120 s, priority 10)
% src: twin/service.py:357
% src: paper/aies2027/kit/FACTS.md:25

\paragraph{The on-robot tier.} The second expert is a local OpenAI-compatible endpoint with a timeout
of 90 seconds. In the twin it is a model on the host and not in the cloud. The fact sheet records it
as gemma-4-26B-A4B, quantised to four bits in GGUF format and served by llama.cpp. While
communications are down the cascade skips the cloud expert and asks only this one.
% src: twin/cascade.py:247-250, 284-288, 310-312
% src: docs/AUTONOMY_PLAN.md:167-168
% src: paper/aies2027/kit/FACTS.md:25 (/archive/unity/tools/brain7.sh)

\paragraph{The compiled tier.} When no expert gives an acceptable answer, the cascade raises the
exception \code{ModelUnavailable} and the brain acts with no model. The scene's offline rules turn the
validated facts into events. Each rule fires when its tokens start to hold and not again while they
keep holding. Smoke on a detector, a fall sensor alert with Margaret lying on the floor, ten minutes
on the floor, and a scream heard from her are examples. The brain then takes the highest-priority
obligation in force that is allowed. A request to the governor names its target from the scene's
table, so \code{p5} asks for \code{restrain\_person} and \code{p3} for \code{drive\_off\_animal}. With
no obligation in force, the default action, chores, stands. The governor and the scene run on the
robot in every tier.
% src: twin/cascade.py:263-264, 295-298
% src: twin/brain.py:508-570 (compiled tier)
% src: twin/scene/margaret_home.erisml:1035-1061 (offline)
% src: docs/AUTONOMY_PLAN.md:169-174

\section{The hash-chained log and the live ladder}\label{sec:home-log}

\paragraph{The log.} Every record the service returns is first appended to an append-only log, one
JSON record per line. Each record carries its sequence number, the hash of the record before it and
a time. The service writes the record as canonical JSON with the audit module's serializer and
stores the SHA-256 of that text as the record's hash. The first record links to a hash of 64 zeros.
Verification walks the chain and checks, for every record, that the canonical text matches its
fields, that the hash matches the text, and that the link matches the previous hash. A third party
can therefore re-verify the chain. In Winfield and Jirotka's terms the log is an ethical black box
\cite{winfield2017ebb}.
% src: twin/service.py:5-7, 49-84 (Chain)

\paragraph{The live ladder.} The escalation ladder of Section~\ref{sec:cont-ladder} is declared in the
scene as a BPMN~2.0 process \cite{omg2011bpmn}. The service exports it at
\code{/process/escalation\_ladder.bpmn} and computes a trace at
\code{/process/escalation\_ladder/trace}. The trace flattens the run since the last reset into the
event stream that erisml-compiler maps onto the process. It also lists where the structure stepped
in, meaning governor refusals, output-gate vetoes and referrals, and lapses. The page
\code{web/ladder.html} draws the process with bpmn-js, polls the trace once a second, marks the
elements taken in this run and the one the robot is at now, and lists the interventions beside the
diagram. The view decides nothing. It shows only what the records say happened. BPMN has been used
before to model and run robot missions \cite{corradini2023bpmn}. Here it is a view of the norms and
not a second source of authority.
% src: twin/service.py:244-264
% src: twin/process_view.py:1-10, 279-352
% src: web/ladder.html:39-49, 63-89
% src: twin/scene/margaret_home.erisml:923-928
